\section{A joint bound for atypical trajectories}
\label{sec:atypical}

The bridge criteria require a bound for many failures at once.  A
single-bridge estimate is insufficient here: the interacting closed measure
does not make its bridges independent.  We obtain a localized exponential
moment by assigning an auxiliary color to each tested bridge, retaining its
spatial confinement, and estimating the resulting partition function.

We use the free-energy increment bound \eqref{eq:6} and the density bound
\eqref{eq:10}.  More explicitly, the two inputs needed in this section are
\[
 F_j-F_{j-1}\le C_v\rho\quad(1\le j\le N),
 \qquad
 \langle n_B\rangle_{J,-W}\le C_v\rho|B|
 \quad(0\le J\le N, 0\le W\le\rho),
\]
where $B$ is any union of coarse cells and the second expectation is in the
$J$-particle equilibrium state with external potential $-W$.  In particular,
the constants in these inputs are fixed before the trajectory cutoff $D$
and the isolation parameter $\zeta$ are chosen.

Recall that $\delta=s^2$, $r=\zeta\rho^{-1/3}$, and the paths in the slab
carry the lift convention specified in Section~\ref{sec:slab}.  For a
nonempty union $U$ of coarse cells, let $J_U^{\mathrm{fail}}$ count the
slots ending in $U$ for which at least one of the following holds: the
lifted path leaves the ball of radius $Ds$ about its end; on the first
time interval of length $r$ it leaves the ball of radius $r/10$ about
its start; or on the last such interval it leaves that ball about its
end.

\begin{proposition}[Localized exponential moment for trajectory failures]
\label{prop:atypical-moment}
Fix $K$ as in the cell construction.  There exist $D_{\min}$, $C_v$, and
$\kappa>0$, depending on the fixed potential and on $K$, such that the
following holds.  Fix $D\ge D_{\min}$ and $0<\zeta\le1$.  For sufficiently
small $\rho>0$ and sufficiently large volume at that density, uniformly
for $\rho/2\le N/V\le2\rho$ and $0<T\le\rho^2$, every nonempty coarse
cell union $U$ satisfies
\begin{equation}
 \log\mathbb E_{\mathrm{cl}}\exp(J_U^{\mathrm{fail}})
 \le
 \left[C_v\sum_{j\ge0}D_j^3e^{-\kappa D_j^2}
       +o_\rho(1)\right]\rho|U|,
 \qquad D_j=2^jD.
 \tag{17}\label{eq:17}
\end{equation}
The constants $C_v,\kappa$ do not depend on $D$ or $\zeta$.  The term
$o_\rho(1)$ is at fixed $v,K,D,\zeta$ and is uniform in the stated
particle numbers, volumes, and temperatures.
\end{proposition}

\subsection{A localized kernel for one marked bridge}

Partition the failure event into disjoint bins.  In bin $j\ge0$, the
lifted maximum distance from the end belongs to $(D_js,2D_js]$.  Bin
$*$ consists of a failure of one of the two short-interval tests with
no macro failure.  Write
\[
 h_j=D_j,\quad b_j=D_j^2\quad(j\ge0),
 \qquad h_*=D,\quad b_*=r.
\]
Here and below a sum over bins includes $*$.  Let
\[
 \Omega_j=\{x\in\Lambda_L:d(x,U)<(2h_j+4)s\}.
\]
If this set is the whole torus, its Laplacian has periodic boundary
conditions; otherwise we use the Dirichlet Laplacian.  Choose a
Lipschitz function $0\le p_j\le1$, equal to one when $d(x,U)\le2h_js$,
vanishing when $d(x,U)\ge(2h_j+2)s$, and satisfying
$|\nabla p_j|\le C/s$.  It has support compactly contained in $\Omega_j$
unless $\Omega_j$ is the whole torus.  The elementary coarse-cell cover
bound gives
\[
 |\Omega_j|\le C h_j^3|U|,
 \qquad
 \Omega_j\subset B_j,\qquad |B_j|\le C h_j^3|U|,
\]
for suitable coarse-cell unions $B_j$.  The constants also work when
the enlarged set covers or winds around the torus.

Let $K_j(E,S)$ be the unnormalized free single-bridge kernel with end
$E\in U$ and failure in bin $j$.  There are absolute constants $c,C>0$
and a fixed $c_1>1$ such that
\begin{equation}
 K_j(E,S)
 \le C e^{-cb_j}p_j(E)p_j(S)
          [e^{c_1\delta\Delta_{\Omega_j}}](E,S).
 \tag{18}\label{eq:18}
\end{equation}
These constants are independent of $D$ and $\zeta$ after the
small-density restrictions below are imposed.

To prove this, fix a lift of $E$ and consider separately each lift of
$S$, with displacement $d$ from that lift of $E$.  Put
$k_t(d)=(4\pi t)^{-3/2}e^{-|d|^2/(4t)}$.  In a macro bin the event
forces $|d|\le2h_js$.  If $|d|>h_js/2$, widening the Gaussian gives
\[
 k_\delta(d)\le C e^{-ch_j^2}k_{c_1\delta}(d).
\]
If $|d|\le h_js/2$, the zero-end bridge fluctuation must have maximum
at least $h_js/2$.  The Brownian bridge maximum tail is at most
$Ce^{-ch_j^2}$, while $k_\delta\le c_1^{3/2}k_{c_1\delta}$.
This proves the same bound for this case.

For bin $*$ one has $|d|\le Ds$.  Once $s\ge40D$, the linear drift
on a time interval of length $r$ is at most $r/40$.  If $B$ is
Brownian motion with generator $\Delta$, the fluctuation is
$B_t-(t/\delta)B_\delta$.  The maximal inequality on $[0,r]$ bounds
the chance that the first term exceeds a fixed positive multiple
of $r$ by $Ce^{-cr}$.  The corresponding estimate for the second
term is $Ce^{-c\delta}$, since exceeding that multiple of $r$
requires $|B_\delta|\ge c\delta$.  We take $r\le\delta/2$.
Time reversal gives the same bound at the other endpoint.  Thus
the bin-$*$ kernel for each retained lift is at most
$Ce^{-cr}k_{c_1\delta}(d)$.

It remains to justify the Dirichlet kernel on the right of
\eqref{eq:18}, rather than a free kernel without localization.
For every retained lift, the straight segment from $E$ to $S$
is within distance $2h_js$ of the chosen lift of $E$.  Every lifted
path at distance less than $s$ from that segment projects into
$\Omega_j$.  A bridge of diffusion coefficient $c_1$ over time
$\delta=s^2$ has probability at least a fixed $q_*>0$ of remaining
in this tube, by scaling its zero-end fluctuation.  Consequently
the killed torus kernel is at least
\[
 q_*\sum_{\substack{\text{retained lifts of }S}}
                  k_{c_1\delta}(d).
\]
The terms are disjoint winding classes of torus paths.  This remains
valid for extents greater than $L$ and for $\Omega_j=\Lambda_L$.
On the support of $K_j$, both $p_j(E)$ and $p_j(S)$ equal one.
Summing the preceding bounds over lifts proves \eqref{eq:18}.

\subsection{The colored operator and its large powers}

Choose $\kappa>0$ sufficiently small, and set
\[
 \eta_j=e^{-\kappa b_j},\qquad M_j=\kappa b_j/\delta.
\]
By increasing the lower restrictions on $D$ and $r$, the prefactor
in \eqref{eq:18}, including the marking factor $e-1$, is bounded by
$\eta_j e^{-\delta M_j}=e^{-2\kappa b_j}$.  The choice of $\kappa$
is independent of $D$ and $\zeta$.

We will represent marked bridges by auxiliary particle colors. The small
coefficient $\eta_j$ limits the enlargement of the map into the colored
space, while the positive gap $M_j$ will pay the free-energy cost of
removing a particle from the original interacting color. Normalizing this
map will introduce a local attractive potential. The density estimate
controls its cost on the original color, and the gap also absorbs it on
the marked colors.

The enlarged one-particle Hilbert space is
\[
 \mathfrak h_e=L^2(\Lambda_L)
                    \oplus\bigoplus_j L^2(\Omega_j),
\]
with the first summand called the original color and the others
indexed by the failure bins.  Define
\[
 v_e f=\bigl(f,(\sqrt{\eta_j}p_jf)_j\bigr),
 \qquad V_e=v_e^{\otimes N}.
\]
The restrictions to the relevant domains in each marked component
are understood.  On the labeled $N$-fold tensor product, define
$\mathbb H$ by its closed quadratic form, separately in each
color assignment.  Original-color particles have their original
kinetic energy and the full pair potential between themselves.
A particle of color $j$ has one-body operator
$-c_1\Delta_{\Omega_j}+M_j$.  All interactions incident to marked
particles are omitted.  The operator and the embedding commute
with, or intertwine, coordinate permutations.  They therefore
restrict to the symmetric $N$-particle spaces.  On those spaces put
\[
 \mathcal T=V_e^*e^{-\delta\mathbb H}V_e.
\]

Indeed, expanding
\[
 e^{J_U^{\mathrm{fail}}}
 =\prod_{i=1}^N\left(1+(e-1)
        \mathbf1_{\{E_i\in U,\ i\text{ fails}\}}\right)
\]
marks subsets of labels.  For each marked subset, decompose its
failure events into bins and drop precisely the interactions
incident to marked labels.  This increases the nonnegative path
weight.  Applying \eqref{eq:18} to each marked path leaves the
original interacting kernel on the other labels.  The sum over
all color assignments is exactly the kernel of $\mathcal T$.
Thus $\mathcal T$ majorizes the tested slab kernel pointwise on
labeled coordinates.  The nonnegative outer kernel, including
its symmetrization projector, gives
\[
 \mathbb E_{\mathrm{cl}}e^{J_U^{\mathrm{fail}}}
 \le Z_N^{-1}\Tr_{\mathrm{sym}}
                     e^{-(\beta-\delta)H}\mathcal T.
\]
For $\beta>\delta$, Schatten H\"older with conjugate exponents
$\beta/(\beta-\delta)$ and $\beta/\delta$ yields
\begin{equation}
 \mathbb E_{\mathrm{cl}}e^{J_U^{\mathrm{fail}}}
 \le
 \left(\frac{\Tr_{\mathrm{sym}}\mathcal T^{\beta/\delta}}
                   {Z_N}\right)^{\delta/\beta}.
 \tag{19}\label{eq:19}
\end{equation}

The embedding has norm larger than one.  We next account for its
effect on the high power in \eqref{eq:19} without losing the spatial
localization.  Define the bounded Lipschitz function
\[
 a(x)=\frac12\log\left(1+\sum_j\eta_jp_j(x)^2\right),
 \qquad A=\sum_{i=1}^N a(x_i).
\]
The same position multiplier acts in every color.  One has
$V_e=e^A U_e$, where $U_e$ is an isometry: at the one-particle
level the squared norm of the color vector is $e^{2a(x)}$.
Consequently
\[
 \mathcal T=U_e^*XU_e,\qquad
 X=e^Ae^{-\delta\mathbb H}e^A.
\]
Eigenvalue min--max for a positive compact compression gives
$\lambda_k(\mathcal T)\le\lambda_k(X)$ for every $k$, including
multiplicities.  Hence the trace of any positive power of
$\mathcal T$ is bounded by that of $X$.

\begin{lemma}[Trace bound for a Lipschitz sandwich]
\label{lem:operator-sandwich}
For the operators just defined, put
\[
 W_0=2a/\delta+c_1|\nabla a|^2,
 \qquad
 \widetilde{\mathbb H}=\mathbb H-\sum_{i=1}^N W_0(x_i).
\]
Then
\begin{equation}
 \Tr X^{\beta/\delta}
       \le\Tr e^{-\beta\widetilde{\mathbb H}}.
 \tag{20}\label{eq:20}
\end{equation}
In fact the analogous inequality holds with every positive power
$q$ on the left and time $q\delta$ on the right.
\end{lemma}

\begin{proof}
We interpolate from the identity to $X$ and bound the rate of change
of the squared norm by the quadratic form of $\widetilde{\mathbb H}$.
Applying the same estimate to exterior powers controls products of
eigenvalues; an integral identity then gives the trace bound for every
positive power.

First, the heat traces below are finite and $X$ is positive compact.
At fixed $N,L,\rho$, the gaps
$M_j$ tend to infinity along the macro bins.  Domain monotonicity
and the torus heat-kernel bound give, for each $t>0$,
\[
 \sum_j e^{-tM_j}\Tr e^{c_1t\Delta_{\Omega_j}}
 \le C V(V^{-1}+t^{-3/2})\sum_j e^{-tM_j}<\infty.
\]
Including the original color, the free labeled $N$-particle trace
is the $N$th power of a finite one-particle trace.  The nonnegative
original-color pair potentials only increase the Hamiltonian in
the form sense.  Min--max therefore gives compact resolvent and
finite positive-time heat traces for $\mathbb H$.  Subtracting
the bounded potential $\sum_iW_0(x_i)$ preserves these properties.
The bounded multiplier $e^A$ also makes $X$ trace class.

For $0\le t\le\delta$ define
\[
 U_t=e^{f(t)A}e^{-t\mathbb H}e^A,
 \qquad f(t)=-1+2t/\delta.
\]
Thus $U_0=1$ and $U_\delta=X$.  Bounded Lipschitz multiplication
preserves the closed form domain of $\mathbb H$.  This includes
the Dirichlet conditions on marked colors and the weighted
sum of color-gap terms.  For any initial vector $x$ and $t>0$,
spectral calculus gives
\[
 e^{-f(t)A}U_tx=e^{-t\mathbb H}e^Ax\in D(\mathbb H).
\]
The vector $y(t)=U_tx$ is norm differentiable for $t>0$, belongs
to the form domain, and satisfies
\[
 y'(t)=\frac{2A}{\delta}y(t)
             -e^{f(t)A}\mathbb H e^{-f(t)A}y(t).
\]
The last expression is well-defined on this evolving vector.
Its real inner product with $y$ can be evaluated using the
closed form between $e^{fA}y$ and $e^{-fA}y$; no assertion that a
Lipschitz multiplier preserves $D(\mathbb H)$ is needed.

In a fixed color component, the diffusion coefficient $c_i$ of
coordinate $i$ is either $1$ or $c_1$.  The product rule gives
almost everywhere
\[
 \operatorname{Re}\left[
  \overline{\nabla_i(e^{fA}y)}\cdot\nabla_i(e^{-fA}y)\right]
 =|\nabla_i y|^2-f^2|\nabla a(x_i)|^2|y|^2.
\]
The cross terms are purely imaginary and all potentials commute
with these multipliers.  Since $|f|\le1$ and $c_i\le c_1$, the
preceding identities imply
\[
 \frac12\frac{\dd}{\dd t}\|U_tx\|^2
 \le-\widetilde{\mathfrak q}[U_tx],
\]
where $\widetilde{\mathfrak q}$ is the form of
$\widetilde{\mathbb H}$.

Let $E_1\le E_2\le\cdots$ be the eigenvalues of
$\widetilde{\mathbb H}$, counting multiplicity.  Apply the same
calculation to the tensor power $U_t^{\otimes k}$, with the sums
of $A$ and $\mathbb H$ over $k$ copies of the enlarged symmetric
many-particle Hilbert space.  Restrict these copies to their
antisymmetric subspace.  In the eigenbasis of
$\widetilde{\mathbb H}$, the smallest eigenvalue of the copy sum
there is $E_1+\cdots+E_k$.  It follows that
\[
 \frac12\frac{\dd}{\dd t}\|U_t^{\otimes k}x\|^2
 \le-(E_1+\cdots+E_k)\|U_t^{\otimes k}x\|^2
\]
for every antisymmetric $x$ and $t>0$.  Integrate from a positive
time and then let that time tend to zero, using $U_0=1$ and
strong continuity.  The norm of the $k$th exterior power of
$U_\delta$ is consequently at most
$e^{-\delta(E_1+\cdots+E_k)}$.

The exterior-power norm is the product of the largest $k$
singular values.  Since $U_\delta=X$ is positive and injective,
write its eigenvalues in decreasing order as $x_i>0$.  We obtain
\[
 \sum_{i=1}^k\log x_i
       \le-\delta\sum_{i=1}^k E_i\qquad(k\ge1).
\]
These partial-sum inequalities imply, for every real $t$,
\[
 \sum_i(\log x_i-t)_+
       \le\sum_i(-\delta E_i-t)_+.
\]
Indeed each side is the supremum over $k\ge0$ of the
corresponding shifted partial sum; for fixed $t$ only finitely
many positive terms occur.  For every $q>0$ the identity
\[
 e^{qu}=\int_{\mathbb R}(u-t)_+q^2e^{qt}\,\dd t
\]
and Tonelli's theorem now give
$\sum_i x_i^q\le\sum_i e^{-q\delta E_i}$.  Taking
$q=\beta/\delta$ proves \eqref{eq:20}.
\end{proof}

\subsection{The partition-function estimate with confined colors}

The preceding lemma has converted the large power of the marked-slab
operator into an ordinary heat trace. It remains to compare that trace
with the original partition function. Spatial confinement is essential
here: it makes the final cost proportional to $|U|$, the volume where
failures are tested, rather than to the whole torus.

Set
\[
 \Sigma=\sum_j\eta_jh_j^3
       =\sum_{j\ge0}D_j^3e^{-\kappa D_j^2}
                            +D^3e^{-\kappa r}.
\]
From the definition of $a$ and Cauchy--Schwarz,
\[
 0\le W_0\le\frac C\delta\sum_j\eta_j\mathbf1_{\Omega_j}.
\]
For example,
\[
 |\nabla a|^2
 =\frac{|\sum_j\eta_jp_j\nabla p_j|^2}
                  {(1+\sum_j\eta_jp_j^2)^2}
 \le\frac C\delta\sum_j\eta_j\mathbf1_{\Omega_j}.
\]
All constants are independent of the number or sizes of the
domains.  Taking $D$ large and then $\rho$ small makes
$\|W_0\|_\infty\le\rho$, because
$\delta\ge K^2\rho^{-1}$ and $r\to\infty$.

For $0\le J\le N$, differentiate the original-color partition
function in the strength of $-W_0$.  The density input, applied
to the covers $B_j$ and to every external potential $-tW_0$ with
$0\le t\le1$, gives
\[
 \log Z_J^{(-W_0)}-\log Z_J
 =\beta\int_0^1\langle n_{W_0}\rangle_{J,-tW_0}\,\dd t
 \le\frac{\beta C_v\rho|U|}\delta\Sigma.
\]
The free-energy increment input also gives
$Z_J/Z_N\le e^{\beta C_v\rho(N-J)}$.

To use these inequalities, decompose the enlarged symmetric
space according to the occupation of each color:
\[
 \operatorname{Sym}^N\mathfrak h_e
 \cong
 \bigoplus_{J+\sum_j k_j=N}
  \operatorname{Sym}^J L^2(\Lambda_L)
       \otimes\bigotimes_j\operatorname{Sym}^{k_j}L^2(\Omega_j).
\]
Each number sector is one tensor product.  There is no
multinomial multiplicity in its trace: the normalized
symmetrization over color assignments is the unitary realizing
this decomposition.  In that sector, the cost
$e^{\beta C_v\rho(N-J)}$ is one factor per marked particle.
On color $j$, subtracting $W_0$ and this cost leaves an operator
bounded below by
\[
 -c_1\Delta_{\Omega_j}
                  +M_j-\|W_0\|_\infty-C_v\rho
 \ge -c_1\Delta_{\Omega_j}+M_j/2.
\]
For this last inequality use
$M_j\ge\kappa b_j\rho/(4K^2)$, choose $D$ sufficiently large
depending on $v,K$, and then make $r$ sufficiently large.
The same choices ensure $\beta M_j/2\ge1$ in every color.

Sum the marked occupation numbers and drop their total-number
constraint.  This bounds the sum by a product of ideal Bose
grand partition functions with gaps $M_j/2$.  In each color its
logarithm is at most
\[
 \sum_{\ell\ge1}\frac{e^{-\ell\beta M_j/2}}\ell
                       \Tr e^{c_1\ell\beta\Delta_{\Omega_j}}
 \le C|\Omega_j|(V^{-1}+T^{3/2})e^{-\beta M_j/2}.
\]
The estimate uses the killed heat-kernel bound by the periodic
one and the geometric series with $\beta M_j/2\ge1$.
The resulting sum over colors is finite, so the passage from
finitely many colors to countably many is justified by
monotone convergence of the nonnegative occupation sums.
Combining this bound with \eqref{eq:19}, the compression
inequality, and \eqref{eq:20}, we obtain the quantitative estimate
\begin{equation*}
 \log\mathbb E_{\mathrm{cl}}e^{J_U^{\mathrm{fail}}}
 \le C_v\rho|U|\Sigma
  +C\delta T|U|(V^{-1}+T^{3/2})
                     \sum_jh_j^3e^{-\beta M_j/2}.
\end{equation*}

This also proves the required uniformity.  Since $\beta\ge\delta$,
the color sum is bounded at fixed $D$ for all sufficiently small
$\rho$, using the summable macro terms
$D_j^3e^{-\kappa D_j^2/2}$ and the micro term
$D^3e^{-\kappa r/2}$.  Moreover $V\ge s^3$,
$\delta\le4K^2/\rho$, and $T\le\rho^2$, so
\[
 \frac{\delta T}{\rho}(V^{-1}+T^{3/2})
 \le C_K(\rho^{3/2}+\rho^3).
\]
The micro contribution $D^3e^{-\kappa r}$ to $\Sigma$ tends
to zero at fixed $D,\zeta$.  Substituting these estimates in the
partition-function bound proves \eqref{eq:17}, with
all errors uniform in the temperature and the thermodynamic
volume, as asserted.
